\section{项目的主要内容和技术路线}

\subsection{主要研究内容}

围绕在视角覆盖不足条件下利用鱼眼相机与雷达几何约束提升 3DGS 室内场景重建质量这一目标，本研究的主要内容拟包括以下四个方面：

\begin{enumerate}
    \item \textbf{鱼眼相机成像建模与标定研究。} 本部分首先研究 OpenCV fisheye 相机模型，明确鱼眼相机在大视场角条件下的投影关系以及焦距、主点和畸变系数等参数化形式。在标定方面，研究内容分为两个层次：其一，利用鱼眼相机采集的多张包含 Charuco 标定板的图像，对鱼眼相机内参进行标定，获得较为准确的焦距、主点和畸变系数等参数；其二，在完成内参标定后，进一步探索更加稳定的位姿恢复算法，以获得更加稳定和准确的相机内外参，并分析鱼眼模型过拟合对参数估计和后续重建质量的影响。
    
    \item \textbf{面向鱼眼相机的 3DGS 重建方法研究。} 传统 3DGS 主要面向透视投影设计，不能直接处理鱼眼图像输入。因此，本部分拟采用将鱼眼图拆分为多张透视图的思路，构建适配传统 3DGS 的场景重建流程。具体而言，研究如何在已知鱼眼相机内参与畸变参数的条件下，通过射线空间采样将单张鱼眼图像转换为若干组虚拟透视图，并根据鱼眼相机外参与虚拟透视相机的相对旋转关系直接计算这些透视图的外参；在此基础上，进一步探索对拆分后的透视图重新进行特征提取、特征匹配与光束法平差优化，以继续提升相机外参精度。最后，将拆分后的透视图接入传统 3DGS 训练管线，并结合更优的高斯增密策略提升场景细节学习能力与最终重建质量。
    
    \item \textbf{雷达点云筛选与几何约束构建方法研究。} 针对雷达点云稀疏、噪声大、离群点多等特点，本部分研究如何获得更加干净、稳定的雷达点云，包括离群点剔除、噪声抑制以及对有效结构点的筛选。在此基础上，进一步探索如何将雷达点云转化为适合 3DGS 优化的几何约束形式，例如深度一致性约束、法线约束约束、结构位置约束等，使雷达信息能够真正参与到 3DGS 的几何优化过程中，从而增强重建结果的结构稳定性并抑制漂浮伪影。
    
    \item \textbf{鱼眼与雷达融合重建系统的验证与实验评估。} 在完成鱼眼相机建模、标定、鱼眼图拆分与 3DGS 适配以及雷达约束构建后，本部分将搭建完整的鱼眼 + 雷达室内场景重建系统，并通过系统性实验进行验证。评估内容包括与其他代表性方法，如 3DGEER \cite{huang20253dgeer}、3DGUT \cite{wu20253dgut} 等工作的量化对比，并统一采用 PSNR、SSIM 和 LPIPS 三项指标进行评价；同时通过消融实验验证鱼眼图拆分策略、透视图外参进一步优化流程、高斯增密策略、雷达点云筛选策略以及不同几何约束设计对最终重建结果的影响。
\end{enumerate}

\subsection{技术路线}

本研究的技术路线可概括为鱼眼相机建模与稳定标定、面向鱼眼的 3DGS 重建、雷达点云约束构建以及系统验证评估四个阶段，其具体流程如下：

\begin{enumerate}
    \item \textbf{数据采集与鱼眼相机建模。} 
    首先完成实验平台搭建，利用鱼眼相机采集室内场景图像，并额外采集多组包含 Charuco 标定板的图像用于相机标定；
    同步采集对应场景的雷达数据。随后基于 OpenCV fisheye 模型建立鱼眼相机成像表达，为后续内参标定、位姿恢复和鱼眼 3DGS 投影建模提供统一的相机模型基础。

    \item \textbf{鱼眼相机内参与位姿恢复。} 
    在标定阶段，首先利用 Charuco 标定板图像完成鱼眼相机内参估计，获得焦距、主点及畸变系数等参数。然后在场景图像上进行特征提取与匹配，并以已标定内参作为先验或约束，通过运动恢复结构与光束法平差优化相机外参及稀疏点云。
    具体实现上，将重点探索以 SuperPoint 为代表的学习式特征提取、以 SuperGlue 或 LightGlue 为代表的高质量匹配方法，并将其接入基于 COLMAP 的增量式 SfM 与光束法平差流程，
    必要时可借鉴 Detector-Free SfM 的相关思路，以提升鱼眼视频帧在弱纹理和强畸变条件下的位姿恢复稳定性。
    该阶段的目标是降低鱼眼相机模型参数与位姿之间的耦合带来的标定结果不准确问题，获得更加稳定、准确的相机内外参。

    \item \textbf{基于射线空间采样的鱼眼图拆分与 3DGS 重建。} 
    在获得可靠的鱼眼相机内外参后，本研究不直接重写整套鱼眼 3DGS 渲染管线，而是首先面向单张已知内参的鱼眼图像，构建若干组给定内参的虚拟透视相机，并将鱼眼图像拆分为多张透视图。
    具体而言，对于虚拟透视图中的任意一个像素，可将其视为一条以鱼眼相机光心为起点的射线；通过对该射线方向施加预设旋转，并结合鱼眼相机的内参和畸变参数，计算该射线方向在原始鱼眼图像上的对应像素位置，从而完成从鱼眼图到透视图的重采样。
    通过合理设置虚拟透视相机的视场角、内参以及旋转角度，可使一张鱼眼图像被拆分为多张能够覆盖原始鱼眼图主要像素区域的透视图。
    对于这些拆分得到的透视图，其外参可以由鱼眼相机外参与虚拟透视相机相对于鱼眼相机的旋转关系直接计算得到，从而获得一组可直接用于传统 3DGS 训练的初始相机参数。
    在此基础上，还可以进一步对拆分后的透视图重新进行特征提取、特征匹配与光束法平差优化，以继续精化透视图外参，提升多视角之间的几何一致性和后续重建质量。
    最后，将这些透视图接入传统 3DGS 训练管线，并进一步引入更优的高斯增密策略，在梯度敏感和细节丰富区域进行更加有效的高斯点增密，以提升场景细节学习能力和最终重建质量。

    \item \textbf{雷达点云辅助初始化与几何约束引入。} 
    在雷达数据侧，首先对原始雷达点云进行质量筛选、噪声抑制和离群点剔除，尽可能提取稳定且有助于表达场景几何结构的有效点。
    在此基础上，探索雷达点云的三种利用方式：
    第一，利用筛选后的雷达点云为高斯点提供更合理的初始空间分布，从而改进 3DGS 中高斯位置的初始化；
    第二，利用逐帧雷达数据进行几何上的相机标定，提前获得较为精确的相机外参初始化，并进一步辅助鱼眼图像的位姿恢复与相机内外参精化；
    第三，根据雷达点云重建场景的粗表面表示，例如尝试采用类似 Marching Cubes 的方法恢复场景表面，并在后续优化中尽量将高斯点约束在该表面附近，将其作为 3DGS 的几何约束，从而辅助提升低纹理区域的跨视角几何一致性。

    \item \textbf{联合优化与系统实验评估。} 
    在鱼眼 3DGS 基线系统中进一步加入雷达几何约束，形成完整的鱼眼 + 雷达联合重建框架。
    最后设计系统性的实验方案，对本方法与 3DGEER、3DGUT 等代表性方法进行对比验证，并统一采用 PSNR、SSIM 和 LPIPS 三项指标进行量化评价；
    同时通过消融实验分析鱼眼图拆分方式、透视图外参进一步优化流程、高斯增密策略、雷达点云筛选方法及不同约束设计对最终重建结果的影响。
\end{enumerate}

\subsection{可行性分析}

\subsubsection{鱼眼建模与标定的可行性}

本研究在鱼眼相机建模与标定方面具有较强可行性。
首先，针对广角和鱼眼相机的几何标定问题，已有综述与比较研究系统梳理了常见相机模型、标定方法及其适用条件，为本研究开展鱼眼相机内参标定提供了较明确的理论与工程参考\cite{huai2024wideanglecalib}。
其次，OpenCV fisheye 模型已经提供了较成熟的参数化表达方式，能够较好描述鱼眼镜头在大视场角条件下的投影关系；
基于 Charuco 标定板的相机标定方法也较为成熟，能够同时利用棋盘格角点与 ArUco 标记信息，在鱼眼图像中获得较稳定的角点检测结果，因此利用多张包含 Charuco 标定板的鱼眼图像完成内参标定，在工程实现上具有明确路径。
进一步地，在获得较可靠内参后，再通过特征提取、特征匹配和光束法平差进行位姿恢复，也符合当前多视图几何的常见流程；
对于鱼眼视频拆分得到的图像帧，其相机外参恢复与优化也可以基于 COLMAP 所采用的增量式运动恢复结构与光束法平差流程实现\cite{schonberger2016sfm}。
同时，SuperPoint、SuperGlue、LightGlue 以及 Detector-Free SfM 等学习式前端与重建流程，已经为弱纹理场景下更稳定的特征提取、匹配与位姿恢复提供了可借鉴的技术基础\cite{detone2018superpoint,sarlin2020superglue,lindenberger2023lightglue,he2024detector}。
因此，该部分在理论基础和工程实现上都具有较好的可操作性。

\subsubsection{鱼眼图拆分与 3DGS 适配的可行性}

本研究提出的将鱼眼图拆分为多张透视图的方案，本质上是一种基于射线空间采样的图像重采样过程，其几何关系清晰、实现链路明确。
首先，对于虚拟透视图上的每个像素，都可以对应到一条以鱼眼相机光心为起点的射线方向，再通过鱼眼相机的内参与畸变模型找到其在原始鱼眼图像上的对应采样位置，因此该过程不依赖额外的学习模块，而主要依赖已经标定好的鱼眼相机参数和预设的虚拟相机旋转角度。
其次，通过设置多组虚拟透视相机，可以较完整地覆盖原始鱼眼图中的有效像素，并将鱼眼观测转化为传统 3DGS 可以直接处理的透视图输入；
这些透视图的外参还可以由鱼眼相机外参与虚拟相机相对姿态直接计算得到，从而较方便地接入现有 3DGS 训练流程。
进一步地，在完成初始拆分后，还可以对透视图重新进行特征提取、特征匹配与光束法平差优化，以继续精化外参并提升几何一致性；
这一步同样可以借助 COLMAP 的增量式 SfM 与光束法平差流程实现\cite{schonberger2016sfm}。
这一思路能够避免从底层重写整个鱼眼 3DGS 渲染框架，显著降低实现复杂度；同时，更优的高斯增密算法已经是当前 3DGS 研究中的常见优化方向，将其接入该基线流程也具有现实可行性。

\subsubsection{雷达辅助初始化与几何约束的可行性}

本研究在雷达辅助初始化与几何约束方面同样具有较强可行性。
首先，雷达数据虽然稀疏且噪声较大，但其提供的几何信息与纹理、光照变化相对解耦，这使其在低纹理区域和几何结构约束不足的场景中具有独特优势，因此即使是较稀疏的点云，也能够为场景主要结构提供几何锚点，用于高斯位置初始化。
其次，利用逐帧雷达数据辅助几何标定，从而为相机外参恢复提供较优初始化，也具有明确的工程意义，因为更好的外参初值有助于降低鱼眼位姿优化中的不稳定性。
进一步地，利用雷达点云重建场景粗表面，并将高斯点约束在表面附近，本质上是把传统几何重建结果转化为 3DGS 的结构先验，这与当前通过几何约束提升 3DGS 一致性的研究方向是相符的。
因此，只要雷达数据经过合理筛选与建模，该部分研究具备较明确的实施基础。

\subsubsection{实验结果评估的可行性}

本研究在实验结果评估方面具有较好可行性。
首先，相关对比方法已经存在，如 3DGEER 和 3DGUT 等面向通用相机或畸变相机的高斯渲染方法，可以作为合理的比较基线，从而支持对本文方法进行横向量化比较。
其次，本研究的方法输出仍然是测试视角下的重建渲染结果，因此可以采用统一的图像质量评价方式开展评估。
在评价指标上，本文将统一采用 PSNR、SSIM 和 LPIPS 三项指标，对不同方法以及不同模块设置下的重建结果进行量化分析，并通过消融实验比较鱼眼图拆分方式、透视图外参进一步优化流程、高斯增密策略以及雷达约束设计对最终结果的影响。
因此，从实验结果对比与量化评价的角度看，本课题具有较强的可执行性。
